
\subsection{\texorpdfstring{$A_\infty$}{A-infinity} 双模同态之间的同伦}
\label{subsec:5-2-3}

本小节定义 $A_\infty$ 双模同态之间的同伦\liftIndex{I185}\liftIndex{I145}，
并陈述本节的主要结果定理 5.2.35。
论证方式与第 4 章的滤过 $A_\infty$ 代数情形相同。
令 $D$ 为滤过 $(C_1,C_0)$-$A_\infty$ 双模，
令 $\mathfrak C_i$ 为 $[0,1]\times C_i$
的模型（$i=0,1$）。

\begin{definition}{5.2.21}\label{def:5-2-21}
称滤过 $(\mathfrak C_1,\mathfrak C_0)$-$A_\infty$
双模 $\mathfrak D$ 为
$[0,1]\times D$ 的一个\emph{模型}\liftIndexRangeStart{I344}，
如果它配有下列 $A_\infty$ 双模同态：
\[
\begin{aligned}
\mathrm{Eval}_{s=s_0}:\mathfrak D&\longrightarrow D\liftIndex{I262}
&&\text{位于 }
(\mathrm{Eval}_{s=s_0},\mathrm{Eval}_{s=s_0})
\text{ 之上},\\
\mathrm{Incl}:D&\longrightarrow\mathfrak D\liftIndex{I142}
&&\text{位于 }(\mathrm{Incl},\mathrm{Incl})
\text{ 之上},
\end{aligned}
\]
其中 $s_0=0,1$，并满足：
\begin{enumerate}[label=\textup{(5.2.22.\arabic*)},leftmargin=5.8em]
\item\label{cond:5-2-22-1}
$\mathrm{Eval}_{s=s_0}\circ\mathrm{Incl}
=\mathrm{id}$；
\item\label{cond:5-2-22-2}
当 $(k_1,k_0)\ne(0,0)$ 时，
$(\mathrm{Eval}_{s=s_0})_{k_1,k_0}
=(\mathrm{Incl})_{k_1,k_0}=0$；
\item\label{cond:5-2-22-3}
$(\mathrm{Eval}_{s=0})_{0,0}
\oplus(\mathrm{Eval}_{s=1})_{0,0}:
\mathfrak D\to D\oplus D$ 是分裂满射；
\item\label{cond:5-2-22-4}
$(\mathrm{Incl})_{0,0}:\overline D
\to\overline{\mathfrak D}$
在 $\overline{\mathfrak n}_{0,0}$ 复形之间
诱导上链同伦等价。
\end{enumerate}
当上述所有对象、评价同态和包含同态均含幺
（相应地，同伦含幺）时，
称 $\mathfrak D$ 为
$[0,1]\times D$ 的\emph{含幺模型}
（相应地，\emph{同伦含幺模型}）。
非滤过版本也以相同方式定义。

对于 $\Lambda_{\mathrm{nov}}$ 上的滤过
$A_\infty$ 双模 $\widetilde D$，
只需取
\[
\widetilde{\mathfrak D}
=\mathfrak D\widehat\otimes_{\Lambda_{0,\mathrm{nov}}}
\Lambda_{\mathrm{nov}}
\]
作为 $[0,1]\times\widetilde D$ 的模型。
注意，对 $\Lambda_{\mathrm{nov}}$ 上的
滤过 $A_\infty$ 双模，条件 (5.2.22.4)
没有意义。
\end{definition}

\begin{theorem}{5.2.23}\label{thm:5-2-23}
对任意滤过 $(C_1,C_0)$-$A_\infty$ 双模 $D$
和 $[0,1]\times C_i$ 的任意模型
$\mathfrak C_i$，均存在
$[0,1]\times D$ 的模型。
若 $C_i$ 与 $D$ 含幺，则存在含幺模型。
非滤过版本也成立。
若 $C_1,C_0,D$ 均有隙，
则可取 $D$ 的有隙模型。
\end{theorem}

证明要用到以下引理。

\begin{lemma}{5.2.24}\label{lem:5-2-24}
$[0,1]\times D$ 的模型是否存在，
不依赖于 $[0,1]\times C_i$（$i=0,1$）
的模型如何选取。
\end{lemma}

\begin{proof}
假设对于某些 $[0,1]\times C_i$ 的模型
$\mathfrak C_i$，
存在 $[0,1]\times D$ 的模型
$\mathfrak D$，它是滤过
$(\mathfrak C_1,\mathfrak C_0)$-$A_\infty$ 双模。
另取 $[0,1]\times C_i$ 的模型
$\mathfrak C'_i$。
对恒等同态应用定理 4.2.34，
可得与 $\mathrm{Eval}$、$\mathrm{Incl}$
可交换的同态
$\mathfrak{ID}_i:\mathfrak C'_i\to\mathfrak C_i$。
于是由引理 5.2.9，
拉回双模
\[
(\mathfrak{ID}_1,\mathfrak{ID}_0)^*
\mathfrak D
\]
是 $[0,1]\times D$ 的模型。
\end{proof}

\begin{proof}[定理 5.2.23 的证明]
由引理 5.2.24，
只须对某组特定的 $\mathfrak C_i$
构造 $[0,1]\times D$ 的模型。
下面分别考虑
$\mathfrak C_i=\operatorname{Poly}([0,1],C_i)$
和 $\mathfrak C_i=C_i^{[0,1]}$
的情形（见第 4.2.1 小节）。
为证明定理本身，考虑
$\mathfrak C_i=C_i^{[0,1]}$ 即已足够；
但明确描述多项式模型的情形也有用。
本证明以 $\mathfrak N_{k_1,k_0}$
表示 $\mathfrak D$ 上的滤过
$A_\infty$ 双模结构，
以 $\mathfrak n_{k_1,k_0}$
表示 $D$ 上的结构。

先考虑
$\mathfrak C_i=\operatorname{Poly}([0,1],C_i)$。
取形如
\[
y(s)=\sum_i y_i(s)T^{\lambda_i},
\qquad
y'(s)=\sum_i y'_i(s)T^{\lambda_i}
\]
的一对元素，其中 $y_i:[0,1]\to D^k[1]=D^{k+1}$
与 $y'_i:[0,1]\to D^{k-1}[1]=D^k$
均为 $s$ 的多项式，
$\lambda_i\geq0$ 且
$\lim_{i\to\infty}\lambda_i=+\infty$。
把所有这样的
$\mathbf y=(y,y')$ 组成的集合记为
$\operatorname{Poly}([0,1],D)^k$\liftIndex{I362}\liftIndexRangeEnd{I344}。

对
$\mathbf x_j^{(i)}
=(x_j^{(i)},x_j'^{(i)})
\in\operatorname{Poly}([0,1],C_i)$
（$i=0,1$），以及
$\mathbf y=(y,y')\in\operatorname{Poly}([0,1],D)$，
定义 $\bigl(\mathfrak N_{k_1,k_0}
(\mathbf x_1^{(1)},\ldots,\mathbf x_{k_1}^{(1)},
\mathbf y,\mathbf x_1^{(0)},\ldots,
\mathbf x_{k_0}^{(0)})\bigr)(s)
=(\widetilde u(s),\widetilde v(s))$，其中
\[
\begin{aligned}
\widetilde u(s)=
\mathfrak n_{k_1,k_0}\bigl(
&x_1^{(1)}(s),\ldots,x_{k_1}^{(1)}(s),\\
&y(s),x_1^{(0)}(s),\ldots,x_{k_0}^{(0)}(s)
\bigr),
\end{aligned}
\]
而当 $(k_1,k_0)\ne(0,0)$ 时，
$\widetilde v(s)$ 按下式定义
\translatornote{原书实体页 293 的中间一项，
即含 $y'(s)$ 的项，右端末个输入
误写为 $x_{k_1}^{(0)}(s)$；
该输入属于右侧的 $C_0$，其个数是 $k_0$，
且前后两项均以 $x_{k_0}^{(0)}(s)$ 结尾。
此处据此改为 $k_0$；见疑误 A-0271。}：
\[
\begin{aligned}
\widetilde v(s)={}&
\sum_{j=1}^{k_1}(-1)^{\epsilon_j}
\mathfrak n_{k_1,k_0}\bigl(
x_1^{(1)}(s),\ldots,x_j'^{(1)}(s),
\ldots,x_{k_1}^{(1)}(s),\\
&\hspace{8em}y(s),x_1^{(0)}(s),
\ldots,x_{k_0}^{(0)}(s)\bigr)\\
&+(-1)^\epsilon
\mathfrak n_{k_1,k_0}\bigl(
x_1^{(1)}(s),\ldots,x_{k_1}^{(1)}(s),
y'(s),\\
&\hspace{8em}x_1^{(0)}(s),\ldots,
x_{k_0}^{(0)}(s)\bigr)\\
&+\sum_{j=1}^{k_0}(-1)^{\epsilon'_j}
\mathfrak n_{k_1,k_0}\bigl(
x_1^{(1)}(s),\ldots,x_{k_1}^{(1)}(s),
y(s),\\
&\hspace{8em}x_1^{(0)}(s),\ldots,
x_j'^{(0)}(s),\ldots,
x_{k_0}^{(0)}(s)\bigr).
\end{aligned}
\]
在 $(k_1,k_0)=(0,0)$ 时，改定义
\[
\widetilde v(s)
=\frac{\partial y'(s)}{\partial s}
-\mathfrak n_{0,0}(y'(s)).
\]
上式的符号指数是
\[
\begin{aligned}
\epsilon_j
&=\sum_{i=1}^{j-1}\deg' x_i^{(1)}(s)+1,\\
\epsilon
&=\sum_{i=1}^{k_1}\deg' x_i^{(1)}(s)+1,\\
\epsilon'_j
&=\sum_{i=1}^{k_1}\deg' x_i^{(1)}(s)
+\deg' y(s)
+\sum_{i=1}^{j-1}\deg' x_i^{(0)}(s)+1.
\end{aligned}
\]
定义位于
$(\mathrm{Eval}_{s=s_0},\mathrm{Eval}_{s=s_0})$
之上的滤过 $A_\infty$ 双模同态
\[
\mathrm{Eval}_{s=s_0}:
\operatorname{Poly}([0,1],D)\longrightarrow D
\]
为
\[
(\mathrm{Eval}_{s=s_0})_{0,0}(y,y')=y(s_0),
\qquad
(\mathrm{Eval}_{s=s_0})_{k_1,k_0}=0
\quad((k_1,k_0)\ne(0,0)).
\]
再定义位于 $(\mathrm{Incl},\mathrm{Incl})$
之上的
$\mathrm{Incl}:D\to\operatorname{Poly}([0,1],D)$
\translatornote{原书实体页 293 在此把
单一的包含同态写成 $\mathrm{Incl}_{s=s_0}$，
但定义 5.2.21 与条件 (5.2.22.1)
只定义与两个评价同态配对的
同一个 $\mathrm{Incl}$。
此处统一记为 $\mathrm{Incl}$；
见疑误 A-0272。}：
\[
(\mathrm{Incl})_{0,0}(v)=(v,0),
\qquad
(\mathrm{Incl})_{k_1,k_0}=0
\quad((k_1,k_0)\ne(0,0)).
\]
按定义检验，
$\operatorname{Poly}([0,1],D)$
是 $[0,1]\times D$ 的模型；
证明直接，从略。

接着考虑 $\mathfrak C_i=C_i^{[0,1]}$。
回顾
$C_i^{[0,1]}=C_i\oplus C_i[-1]\oplus C_i$
（见第 4.2.1 小节），相应地置
\[
\mathfrak D=D\oplus D[-1]\oplus D.
\]
定义 $\mathfrak I_0,\mathfrak I_1:D\to\mathfrak D$
及 $\mathfrak J:D\to\mathfrak D$：
\[
\mathfrak I_0(u)=(u,0,0),\qquad
\mathfrak I_1(w)=(0,0,w),\qquad
\mathfrak J(v)=(0,v,0).
\]
$\mathfrak I_0,\mathfrak I_1$ 保持次数，
$\mathfrak J$ 的次数为 $+1$。
以 $\mathfrak N$ 记将要定义的
$\mathfrak D$ 上的滤过 $A_\infty$ 双模结构。
对 $u,v,w\in D$，先置下列值
\translatornote{原书实体页 293 的式 (5.2.25)
把前两行符号指数都写作 $\deg' v$，
但这两行的输入分别是 $u,w$；
与第 4.2.1 小节同型的区间模型微分
及各行的次数型别相对照，
应分别为 $\deg' u,\deg' w$。
此处据此修正；见疑误 A-0273。}：
\begin{equation}
\begin{aligned}
\mathfrak N_{0,0}(\mathfrak I_0(u))
&=\mathfrak I_0(\mathfrak n_{0,0}u)
  +(-1)^{\deg' u}\mathfrak J(u),\\
\mathfrak N_{0,0}(\mathfrak I_1(w))
&=\mathfrak I_1(\mathfrak n_{0,0}w)
  -(-1)^{\deg' w}\mathfrak J(w),\\
\mathfrak N_{0,0}(\mathfrak J(v))
&=\mathfrak J(\mathfrak n_{0,0}v).
\end{aligned}
\tag{5.2.25}\label{eq:5-2-25}
\end{equation}
再置
\begin{equation}
\begin{aligned}
(\mathrm{Eval}_{s=0})_1(u,v,w)&=u,
&(\mathrm{Eval}_{s=1})_1(u,v,w)&=w,\\
(\mathrm{Incl})_1(v)
&=\mathfrak I_0(v)+\mathfrak I_1(v)
=(v,0,v).
\end{aligned}
\tag{5.2.26}\label{eq:5-2-26}
\end{equation}

如下定义 $\mathfrak D$ 上的滤过
$A_\infty$ 双模结构。
取 $\mathbf x^{(i)},\mathbf z^{(i)}
\in B C_i[1]$ 及 $y^{(i)}\in C_i[1]$。
对 $u,v,w\in D$，置
\[
\begin{aligned}
&\mathfrak N\bigl(
\mathfrak I_0(\mathbf x^{(1)}),
\mathfrak I_0(u),
\mathfrak I_0(\mathbf x^{(0)})\bigr)\\
&\qquad=\mathfrak I_0\bigl(
\mathfrak n(\mathbf x^{(1)},u,\mathbf x^{(0)})\bigr),\\
&\mathfrak N\bigl(
\mathfrak I_1(\mathbf z^{(1)}),
\mathfrak I_1(w),
\mathfrak I_1(\mathbf z^{(0)})\bigr)\\
&\qquad=\mathfrak I_1\bigl(
\mathfrak n(\mathbf z^{(1)},w,\mathbf z^{(0)})\bigr).
\end{aligned}
\]
混合输入时，定义
\[
\begin{aligned}
&\mathfrak N\bigl(
\mathfrak I_0(\mathbf x^{(1)}),
\mathfrak I_0(u),
\mathfrak I_0(\mathbf x^{(0)}),
\mathfrak J(y^{(0)}),
\mathfrak I_1(\mathbf z^{(0)})\bigr)\\
&\qquad=(-1)^{\deg'\mathbf z^{(0)}}
\mathfrak J\bigl(
\mathfrak n(\mathbf x^{(1)},u,
\mathbf x^{(0)},y^{(0)},\mathbf z^{(0)})\bigr),\\[0.5ex]
&\mathfrak N\bigl(
\mathfrak I_0(\mathbf x^{(1)}),
\mathfrak J(v),
\mathfrak I_1(\mathbf z^{(0)})\bigr)\\
&\qquad=(-1)^{\deg'\mathbf z^{(0)}}
\mathfrak J\bigl(
\mathfrak n(\mathbf x^{(1)},v,\mathbf z^{(0)})\bigr),\\[0.5ex]
&\mathfrak N\bigl(
\mathfrak I_0(\mathbf x^{(1)}),
\mathfrak J(y^{(1)}),
\mathfrak I_1(\mathbf z^{(1)}),
\mathfrak I_1(w),
\mathfrak I_1(\mathbf z^{(0)})\bigr)\\
&\qquad=(-1)^{
\deg'\mathbf z^{(1)}+\deg' w+\deg'\mathbf z^{(0)}}
\mathfrak J\bigl(
\mathfrak n(\mathbf x^{(1)},y^{(1)},
\mathbf z^{(1)},w,\mathbf z^{(0)})\bigr).
\end{aligned}
\]
其他情形一律置 $\mathfrak N=0$。
与引理 4.2.25 的证明相同，
可以证明这定义滤过 $A_\infty$ 双模结构。
非滤过以及／或者含幺、同伦含幺的版本类似。
定理 5.2.23 得证。
\end{proof}

下面的定理 5.2.27 是定理 4.2.34
在双模同态情形下的对应结果，
对研究 $A_\infty$ 双模同态之间的同伦
有基础作用。
设 $D$、$D'$ 分别是滤过
$(C_1,C_0)$-、$(C'_1,C'_0)$-$A_\infty$ 双模；
令 $\mathfrak D$、$\mathfrak D'$
分别是 $[0,1]\times D$、
$[0,1]\times D'$ 的模型，
它们分别是滤过
$(\mathfrak C_1,\mathfrak C_0)$-、
$(\mathfrak C'_1,\mathfrak C'_0)$-$A_\infty$ 双模。
设 $\mathfrak f^{(i)}:C_i\to C'_i$
与 $\mathfrak F^{(i)}:\mathfrak C_i\to\mathfrak C'_i$
为滤过 $A_\infty$ 同态，满足
\[
\begin{aligned}
\mathrm{Eval}_{s=s_0}\circ\mathfrak F^{(i)}
&=\mathfrak f^{(i)}\circ\mathrm{Eval}_{s=s_0},\\
\mathrm{Incl}\circ\mathfrak f^{(i)}
&=\mathfrak F^{(i)}\circ\mathrm{Incl}.
\end{aligned}
\]
又设 $\varphi:D\to D'$ 是位于
$(\mathfrak f^{(1)},\mathfrak f^{(0)})$
之上的滤过 $A_\infty$ 双模同态。
假定上述所有对象和同态均 $G$-有隙。

\begin{theorem}{5.2.27}\label{thm:5-2-27}
存在位于
$(\mathfrak F^{(1)},\mathfrak F^{(0)})$
之上的滤过 $A_\infty$ 双模同态
\[
\varphi^{\times[0,1]}:\mathfrak D
\longrightarrow\mathfrak D',
\]
使对 $s_0=0,1$ 均有
\[
\mathrm{Eval}_{s=s_0}\circ
\varphi^{\times[0,1]}
=\varphi\circ\mathrm{Eval}_{s=s_0},
\qquad
\mathrm{Incl}\circ\varphi
=\varphi^{\times[0,1]}\circ\mathrm{Incl}.
\]
非滤过或弱滤过的版本，
以及／或者含幺或同伦含幺的版本，
也都成立。
\end{theorem}

定理 5.2.27 的证明与定理 4.2.34 的证明平行，
也类似将在第 5.2.5--5.2.7 小节证明的
Whitehead 定理双模版本，即定理 5.2.35。
因此这里略去定理 5.2.27 的一般证明，
只为弱滤过 $A_\infty$ 双模同态情形
证明下面的引理：定理 5.2.35
并不处理这种情形，我们也不知道它在此情形下
应如何正确表述。

\begin{lemma}{5.2.28}\label{lem:5-2-28}
定理 5.2.27 的滤过版本蕴涵相应的弱滤过版本。
具体地，若 $\varphi$ 弱滤过，
则可使 $\varphi^{\times[0,1]}$
成为弱滤过 $A_\infty$ 双模同态；
若 $\varphi$ 的能量损失为 $c$，
则 $\varphi^{\times[0,1]}$
的能量损失也是 $c$。
\end{lemma}

\begin{proof}
假定 $\mathfrak f^{(i)}$、
$\mathfrak F^{(i)}$ 等均 $G$-有隙，
并且 $\widetilde\varphi:
\widetilde D\to\widetilde D'$
对某个满足条件 (5.2.4) 的
$G$-集合 $G'$ 有隙。
令
\[
\lambda_0
=\min\{\lambda'_i\mid
(\lambda'_i,\mu'_i)\in G'\},
\qquad
c=\max\{-\lambda_0,0\},
\]
并置
\[
G'+c
=\{(\lambda'_i+c,\mu'_i)
\mid(\lambda'_i,\mu'_i)\in G'\}.
\]
于是 $G'+c$ 是满足 (5.2.4) 的
$G$-集合，且
$G'+c\subseteq\mathbb R_{\geq0}\times2\mathbb Z$。
若
\[
\widetilde\varphi_{k_1,k_0}
=\sum_i T^{\lambda'_i}e^{\mu'_i/2}
\varphi_{k_1,k_0,i},
\]
则
\[
T^c\widetilde\varphi_{k_1,k_0}
=\sum_i T^{c+\lambda'_i}e^{\mu'_i/2}
\varphi_{k_1,k_0,i}.
\]
可将 $T^c\widetilde\varphi_{k_1,k_0}$
视为位于
$(\mathfrak f^{(1)},\mathfrak f^{(0)})$
之上的 $(G'+c)$-有隙滤过
$A_\infty$ 双模同态 $D\to D'$。
于是对它应用定理 5.2.27 的滤过版本，
得到
$T^c\widetilde\varphi^{\times[0,1]}$。
记 $\langle G'+c\rangle$
为由 $G'+c$ 生成的幺半群；
它满足条件 3.1.6。
再以相同方式乘以 $T^{-c}$，
便得到满足所需性质的
$\langle G'+c\rangle-c$-有隙
弱滤过 $A_\infty$ 双模同态
\[
\widetilde\varphi^{\times[0,1]}:
\widetilde{\mathfrak D}
\longrightarrow\widetilde{\mathfrak D'}.
\]
它的能量损失为 $c$。
这里
\[
\langle G'+c\rangle-c
=\{(\lambda'_i-c,\mu'_i)\mid
(\lambda'_i,\mu'_i)\in\langle G'+c\rangle\}
\]
是一个 $\langle G'+c\rangle$-集合。
\end{proof}

一旦把 $G'$ 扩充为
$G''=\langle G'+c\rangle-c$，
就不必再扩充，因为
$G''=\langle G''+c\rangle-c$。

\begin{definition}{5.2.29}\label{def:5-2-29}
设 $\varphi,\psi:D\to D'$
分别是位于
$(\mathfrak f^{(1)},\mathfrak f^{(0)})$、
$(\mathfrak g^{(1)},\mathfrak g^{(0)})$
之上的滤过 $A_\infty$ 双模同态，
其中 $\mathfrak f^{(i)},\mathfrak g^{(i)}
:C_i\to C'_i$ 是滤过 $A_\infty$ 同态。
若下列条件 (5.2.30) 成立，
称 $\varphi$ \emph{同伦于} $\psi$\liftIndex{I324}，
记作 $\varphi\sim\psi$：
\begin{equation}
\begin{minipage}{0.86\linewidth}
存在 $[0,1]\times D'$ 的模型
$\mathfrak D'$ 与
$[0,1]\times C'_i$ 的模型
$\mathfrak C'_i$，
以及 $\mathfrak f^{(i)}$ 与
$\mathfrak g^{(i)}$ 之间的同伦
$\mathfrak F^{(i)}:C_i\to\mathfrak C'_i$。
并且存在位于
$(\mathfrak F^{(1)},\mathfrak F^{(0)})$
之上的滤过 $A_\infty$ 双模同态
$\Phi:D\to\mathfrak D'$，
满足
$\mathrm{Eval}_{s=0}\circ\Phi=\varphi$
与
$\mathrm{Eval}_{s=1}\circ\Phi=\psi$。
\end{minipage}
\tag{5.2.30}\label{eq:5-2-30}
\end{equation}
称 $\Phi$ 为从 $\varphi$ 到 $\psi$
的一个\emph{同伦}。

若对于给定的
$\mathfrak C'_i,\mathfrak D',
\mathfrak F^{(i)}$
可取上述 $\Phi$，记作
$\varphi\sim_{(\mathfrak C'_i,\mathfrak D',
\mathfrak F^{(i)})}\psi$。
类似地定义只固定其中部分数据时的记号
$\sim_{(\mathfrak C'_i,\mathfrak F^{(i)})}$、
$\sim_{(\mathfrak C'_i,\mathfrak D')}$、
$\sim_{(\mathfrak C'_i)}$。
若 $\mathfrak F^{(i)}$ 严格，
且 $\varphi\sim_{\mathfrak F^{(i)}}\psi$，
也记作 $\varphi\sim_s\psi$。

非滤过或弱滤过的版本，
以及／或者含幺或同伦含幺的版本，
均以相同方式定义。
同伦含幺版本定义为
$\varphi^+,\psi^+$ 的含幺版本
（记号见定义 3.7.12）。
沿用第 4.2.3 小节滤过
$A_\infty$ 代数同态的记号
$\sim_u,\sim_{hu},\sim_{su}$ 等。
\end{definition}

\begin{lemma}{5.2.31}\label{lem:5-2-31}
关系 $\sim_{(\mathfrak C'_i,\mathfrak D')}$
不依赖于 $\mathfrak C'_i,\mathfrak D'$
的选择；
$\sim_{(\mathfrak C'_i)}$
不依赖于 $\mathfrak C'_i$ 的选择；
$\sim_{(\mathfrak C'_i,\mathfrak D',
\mathfrak F^{(i)})}$
不依赖于 $\mathfrak D'$ 的选择。
\end{lemma}

\begin{proof}
假设
$\varphi\sim_{(\mathfrak C'_i,\mathfrak D')}\psi$，
并取同伦 $\Phi$，
它位于
$(\mathfrak F^{(1)},\mathfrak F^{(0)})$
之上。
另选 $\mathfrak C''_i,\mathfrak D''$。
由定理 4.2.34 与定理 5.2.27，
存在滤过 $A_\infty$ 同态
$\mathfrak{ID}_i:\mathfrak C'_i\to\mathfrak C''_i$
及位于
$(\mathfrak{ID}_1,\mathfrak{ID}_0)$
之上的滤过 $A_\infty$ 双模同态
$\mathfrak{ID}:\mathfrak D'\to\mathfrak D''$，
使如下端点评价相容
\translatornote{原书实体页 296 在两处均把右端写成
“identity”；但复合映射的定义域分别是
$\mathfrak C'_i,\mathfrak D'$，
陪域分别是 $C'_i,D'$，
不是这些对象上的恒等映射。
此处把“端点上为恒等”的含义写成
两个评价映射相容式；
见疑误 A-0275。}：
\[
\mathrm{Eval}^{\prime\prime}_{s=s_0}
\circ\mathfrak{ID}_i
=\mathrm{Eval}^{\prime}_{s=s_0},
\qquad
\mathrm{Eval}^{\prime\prime}_{s=s_0}
\circ\mathfrak{ID}
=\mathrm{Eval}^{\prime}_{s=s_0}
\quad(s_0=0,1).
\]
于是 $\mathfrak{ID}\circ\Phi:D\to\mathfrak D''$
给出 $\varphi,\psi$ 之间的同伦，
从而得证第一项。
第二、三项同理。
\end{proof}

\begin{remark}{5.2.32}\label{rem:5-2-32}
关系
$\sim_{(\mathfrak C'_i,\mathfrak D',
\mathfrak F^{(i)})}$
一般似乎依赖于
$\mathfrak F^{(i)}$ 的选择；
但在“相对于边界的同伦”\liftIndex{I339}
意义下改变 $\mathfrak F^{(i)}$ 时，
它是不变的。
本书不使用这一断言，故不予证明。
描述这一情形（以及命题 5.2.43 的情形）
的恰当语言似乎是
\emph{纤维化范畴}（见 \textbf{[Grot71]}）。
\end{remark}

以下不再标明
$\sim_{(\mathfrak C'_i,\mathfrak D')}$、
$\sim_{(\mathfrak C'_i)}$
的下标，而统一写作 $\sim$。
同时以 $\sim_{\mathfrak F^{(i)}}$
代替
$\sim_{(\mathfrak C'_i,\mathfrak F^{(i)})}$、
$\sim_{(\mathfrak C'_i,\mathfrak D',
\mathfrak F^{(i)})}$。

\begin{lemma}{5.2.33}\label{lem:5-2-33}
滤过 $A_\infty$ 双模同态之间的关系 $\sim$
是等价关系，且与同态的复合相容。
\end{lemma}

证明与命题 4.2.37 和引理 4.2.43 类似，
这里从略。
引理 5.2.31 与 5.2.33 的
非滤过或弱滤过版本，
以及／或者含幺或同伦含幺的版本，
也都成立。

\begin{definition}{5.2.34}\label{def:5-2-34}
设 $\varphi:D\to D'$ 为位于
$(\mathfrak f^{(1)},\mathfrak f^{(0)})$
之上的滤过 $A_\infty$ 双模同态，
其中 $\mathfrak f^{(i)}:C_i\to C'_i$
是滤过 $A_\infty$ 同态。
若存在 $\psi:D'\to D$，
使得 $\psi\circ\varphi$
同伦于恒等同态，
就称 $\varphi$ 为\emph{同伦等价}\liftIndex{I186}。
\translatornote{原书实体页 296 的定义
只写 $\psi\circ\varphi$ 同伦于恒等同态，
未写另一方向的
$\varphi\circ\psi$ 条件；
这与第 4.2.42 定义的双边同伦等价
以及下文“同伦逆”的通常读法不完全一致。
此处保持原定义的量词范围，
不擅自补入未写出的条件；
见疑误 A-0274。}
若存在同伦等价
$\varphi:D\to D'$，
称 $D$ 与 $D'$ \emph{同伦等价}。
非滤过或弱滤过的版本，
以及／或者含幺或同伦含幺的版本，
均以相同方式定义。
若 $\mathfrak f^{(i)}$ 严格，
且所涉同伦位于严格滤过
$A_\infty$ 同态之上，
则称 $\varphi$ 为\emph{严格同伦等价}\liftIndex{I187}。
\end{definition}

注意，若位于
$(\mathfrak f^{(1)},\mathfrak f^{(0)})$
之上的
$\varphi:D\to D'$ 是同伦等价，
则 $\mathfrak f^{(i)}$ 必为同伦等价。

现在陈述本节的主要结果，
即滤过 $A_\infty$ 双模同态的
Whitehead 定理版本\liftIndex{I191}。

\begin{theorem}{5.2.35}\label{thm:5-2-35}
设 $\varphi:D\to D'$ 是位于
$(\mathfrak f^{(1)},\mathfrak f^{(0)})$
之上的有隙滤过 $A_\infty$ 双模同态，
其中 $\mathfrak f^{(i)}:C_i\to C'_i$
是同伦等价，
并以 $(\mathfrak f^{(i)})^{-1}$
表示它们的同伦逆。
假设 $\overline\varphi:
\overline D\to\overline D'$
在 $\overline{\mathfrak n}_{0,0}$
上同调上诱导同构。
那么 $\varphi$ 是同伦等价。
可取位于
$\bigl((\mathfrak f^{(1)})^{-1},
(\mathfrak f^{(0)})^{-1}\bigr)$
之上的 $\varphi$ 的同伦逆
$\varphi^{-1}$。
而且，从恒等同态到
$\varphi^{-1}\circ\varphi$ 的同伦，
可位于任意选定的、从恒等同态到
$(\mathfrak f^{(i)})^{-1}\circ
\mathfrak f^{(i)}$ 的同伦之上。
非滤过以及／或者含幺、同伦含幺
以及／或者严格的版本也成立。
\end{theorem}

证明仍按定理 4.2.34 与 4.2.45 的思路，
分别对数目滤过和能量滤过作两次归纳。
定理 5.2.35 将在第 5.2.5--5.2.7 小节证明。

\begin{remark}{5.2.36}\label{rem:5-2-36}
对于 $\Lambda_{\mathrm{nov}}$ 上的
弱滤过 $A_\infty$ 双模同态
（定义 5.2.1），
我们不知道定理 5.2.35 的对应结论，
因为尚不知如何处理能量滤过。
并且在弱滤过情形下，
“$\overline\varphi:\overline D\to
\overline D'$ 在
$\overline{\mathfrak n}_{0,0}$ 上同调上诱导同构”
这一假设本身就没有意义，
因为此时 $\varphi$ 不诱导
$\overline\varphi$。

另一方面，若
$\mathfrak f^{(1)},\mathfrak f^{(0)}$
均严格，也许可以不经能量滤过，
直接使用数目滤过（见子引理 5.2.20）。
但由此所得的“同伦逆”
是否弱滤过，我们并无把握。
处理定理 5.2.27 弱滤过版本的办法
（引理 5.2.28）
对定理 5.2.35 并不适用。
与定理 5.2.35 的弱滤过版本有关的困难，
或许与将在第 6 章证明的定理 J 有关。
\end{remark}
